% GENERATED FILE. Edit canonical lessons, then run npm run generate:books.
% canonical-source-hash: fnv1a64:48011d4733f622c1
% canonical-lessons: GU-C493-pedhi, GU-C493-purvaj, GU-C493-bujhurg, GU-C493-parivar, GU-C493-santan, GU-R493-review-487-493

\chapter{Many Generations}
\label{ch:gu-many-generations}
\hlchaptermodality{493}

\begin{chapteropening}
\textbf{By the end of this chapter:} \emph{I can talk about my family's generations.}

Complete the last lesson of chapter 493: I can talk about my family's generations.
\end{chapteropening}

\section[\texorpdfstring{\gu{પેઢી} (peḍhī)}{પેઢી (peḍhī)}]{\gu{પેઢી} (peḍhī) — a generation}

\label{lesson:GU-C493-pedhi}



\begin{quote}
Before the new one: say the Gujarati for a grandson, then the Gujarati for a cousin.
\end{quote}

\subsection*{You'll want to know: \gu{પેઢી}}
\textbf{\gu{પેઢી}} — \emph{peḍhī} — \textquotedblleft{}a generation\textquotedblright{}.

\begin{quote}
\gu{અમારા} \gu{કુટુંબની} \gu{ચાર} \gu{પેઢી} \gu{આ} \gu{ગામમાં} \gu{રહી} \gu{છે}.
\end{quote}

\begin{quote}
\emph{amārā kuṭumbnī cār peḍhī ā gāmmā̃ rahī che.}
\end{quote}

\begin{quote}
Four generations of our family have lived in this village.
\end{quote}

A feminine noun.

\begin{cousinweb}
\textbf{\gu{પેઢી}} (\emph{peḍhī}) is inherited from Sanskrit \emph{pīṭha}, \textquotedblleft{}stool, bench\textquotedblright{}; CDIAL lists the Gujarati word as \textquotedblleft{}step of a ladder\textquotedblright{} (CDIAL 8222). The sense \textquotedblleft{}a generation\textquotedblright{}, one step in a family line, is probably a later use of the same word.
\end{cousinweb}

\begin{grammarlens}[title={What you've built}]
One more word for describing what you lived and felt.
\end{grammarlens}

\subsection*{Your turn}
\begin{itemize}
\raggedright
  \item \emph{Say it:} \emph{peḍhī}
  \item \emph{Say it:} \emph{peḍhī}, once more
  \item \emph{Say it:} \emph{amārā kuṭumbnī cār peḍhī ā gāmmā̃ rahī che.} — Four generations of our family have lived in this village.
  \item \emph{Recall:} say the Gujarati for a grandson, then the Gujarati for a cousin, then say \emph{peḍhī} again
\end{itemize}

\begin{tcolorbox}[breakable,colback=teal!4,colframe=teal!35!black,title={Before you move on}]
What does \gu{પેઢી} mean? (\textquotedblleft{}A generation\textquotedblright{}.) Say it once more.
\end{tcolorbox}
\section[\texorpdfstring{\gu{પૂર્વજ} (pūrvaj)}{પૂર્વજ (pūrvaj)}]{\gu{પૂર્વજ} (pūrvaj) — an ancestor}

\label{lesson:GU-C493-purvaj}



\begin{quote}
Before the new one: say the Gujarati for a cousin, then the Gujarati for a generation.
\end{quote}

\subsection*{You'll want to know: \gu{પૂર્વજ}}
\textbf{\gu{પૂર્વજ}} — \emph{pūrvaj} — \textquotedblleft{}an ancestor\textquotedblright{}.

\begin{quote}
\gu{અમારા} \gu{પૂર્વજો} \gu{ઘણા} \gu{મહેનતુ} \gu{હતા}.
\end{quote}

\begin{quote}
\emph{amārā pūrvajo ghaṇā mahenatu hatā.}
\end{quote}

\begin{quote}
Our ancestors were very hard-working.
\end{quote}

A masculine noun.

\begin{cousinweb}
\textbf{\gu{પૂર્વજ}} (\emph{pūrvaj}) is a tatsama, Sanskrit \emph{pūrvaja}, \textquotedblleft{}born before, former\textquotedblright{}, as Monier-Williams records it.
\end{cousinweb}

\begin{grammarlens}[title={What you've built}]
One more word for describing what you lived and felt.
\end{grammarlens}

\subsection*{Your turn}
\begin{itemize}
\raggedright
  \item \emph{Say it:} \emph{pūrvaj}
  \item \emph{Say it:} \emph{pūrvaj}, once more
  \item \emph{Say it:} \emph{amārā pūrvajo ghaṇā mahenatu hatā.} — Our ancestors were very hard-working.
  \item \emph{Recall:} say the Gujarati for a cousin, then the Gujarati for a generation, then say \emph{pūrvaj} again
\end{itemize}

\begin{tcolorbox}[breakable,colback=teal!4,colframe=teal!35!black,title={Before you move on}]
What does \gu{પૂર્વજ} mean? (\textquotedblleft{}An ancestor\textquotedblright{}.) Say it once more.
\end{tcolorbox}
\section[\texorpdfstring{\gu{બુઝુર્ગ} (bujhurg)}{બુઝુર્ગ (bujhurg)}]{\gu{બુઝુર્ગ} (bujhurg) — a respected elder}

\label{lesson:GU-C493-bujhurg}



\begin{quote}
Before the new one: say the Gujarati for a generation, then the Gujarati for an ancestor.
\end{quote}

\subsection*{You'll want to know: \gu{બુઝુર્ગ}}
\textbf{\gu{બુઝુર્ગ}} — \emph{bujhurg} — \textquotedblleft{}a respected elder\textquotedblright{}.

\begin{quote}
\gu{ગામના} \gu{બુઝુર્ગો} \gu{સાંજે} \gu{મંદિર} \gu{પાસે} \gu{બેસે} \gu{છે}.
\end{quote}

\begin{quote}
\emph{gāmnā bujhurgo sā̃je mandir pāse bese che.}
\end{quote}

\begin{quote}
In the evenings the village elders sit near the temple.
\end{quote}

A masculine or feminine noun, by the person meant.

\begin{cousinweb}
\textbf{\gu{બુઝુર્ગ}} (\emph{bujhurg}) is usually traced to Persian \emph{buzurg}, \textquotedblleft{}great, elder\textquotedblright{}.
\end{cousinweb}

\begin{grammarlens}[title={What you've built}]
One more word for describing what you lived and felt.
\end{grammarlens}

\subsection*{Your turn}
\begin{itemize}
\raggedright
  \item \emph{Say it:} \emph{bujhurg}
  \item \emph{Say it:} \emph{bujhurg}, once more
  \item \emph{Say it:} \emph{gāmnā bujhurgo sā̃je mandir pāse bese che.} — In the evenings the village elders sit near the temple.
  \item \emph{Recall:} say the Gujarati for a generation, then the Gujarati for an ancestor, then say \emph{bujhurg} again
\end{itemize}

\begin{tcolorbox}[breakable,colback=teal!4,colframe=teal!35!black,title={Before you move on}]
What does \gu{બુઝુર્ગ} mean? (\textquotedblleft{}A respected elder\textquotedblright{}.) Say it once more.
\end{tcolorbox}
\section[\texorpdfstring{\gu{પરિવાર} (parivār)}{પરિવાર (parivār)}]{\gu{પરિવાર} (parivār) — a household; family members}

\label{lesson:GU-C493-parivar}



\begin{quote}
Before the new one: say the Gujarati for an ancestor, then the Gujarati for a respected elder.
\end{quote}

\subsection*{You'll want to know: \gu{પરિવાર}}
\textbf{\gu{પરિવાર}} — \emph{parivār} — \textquotedblleft{}a household; family members\textquotedblright{}.

\begin{quote}
\gu{અમારા} \gu{પરિવારમાં} \gu{દસ} \gu{માણસ} \gu{છે}.
\end{quote}

\begin{quote}
\emph{amārā parivārmā̃ das māṇas che.}
\end{quote}

\begin{quote}
There are ten people in our household.
\end{quote}

A masculine noun.

\begin{cousinweb}
\textbf{\gu{પરિવાર}} (\emph{parivār}) is a tatsama, Sanskrit \emph{parivāra}, \textquotedblleft{}a cover\textquotedblright{} and \textquotedblleft{}the people around someone, a retinue\textquotedblright{}, as Monier-Williams records it.
\end{cousinweb}

\begin{grammarlens}[title={What you've built}]
One more word for describing what you lived and felt.
\end{grammarlens}

\subsection*{Your turn}
\begin{itemize}
\raggedright
  \item \emph{Say it:} \emph{parivār}
  \item \emph{Say it:} \emph{parivār}, once more
  \item \emph{Say it:} \emph{amārā parivārmā̃ das māṇas che.} — There are ten people in our household.
  \item \emph{Recall:} say the Gujarati for an ancestor, then the Gujarati for a respected elder, then say \emph{parivār} again
\end{itemize}

\begin{tcolorbox}[breakable,colback=teal!4,colframe=teal!35!black,title={Before you move on}]
What does \gu{પરિવાર} mean? (\textquotedblleft{}A household; family members\textquotedblright{}.) Say it once more.
\end{tcolorbox}
\section[\texorpdfstring{\gu{સંતાન} (santān)}{સંતાન (santān)}]{\gu{સંતાન} (santān) — offspring; a child}

\label{lesson:GU-C493-santan}



\begin{quote}
Before the new one: say the Gujarati for a respected elder, then the Gujarati for a household.
\end{quote}

\subsection*{You'll want to know: \gu{સંતાન}}
\textbf{\gu{સંતાન}} — \emph{santān} — \textquotedblleft{}offspring; a child\textquotedblright{}.

\begin{quote}
\gu{તેમને} \gu{બે} \gu{સંતાન} \gu{છે}, \gu{એક} \gu{દીકરો} \gu{અને} \gu{એક} \gu{દીકરી}.
\end{quote}

\begin{quote}
\emph{temne be santān che, ek dīkro ane ek dīkrī.}
\end{quote}

\begin{quote}
They have two children, a son and a daughter.
\end{quote}

A neuter noun.

\begin{cousinweb}
\textbf{\gu{સંતાન}} (\emph{santān}) is a tatsama, Sanskrit \emph{saṃtāna}, \textquotedblleft{}continued succession\textquotedblright{}, as Monier-Williams records it: the line that carries a family on.
\end{cousinweb}

\begin{grammarlens}[title={What you've built}]
One more word for describing what you lived and felt.
\end{grammarlens}

\subsection*{Your turn}
\begin{itemize}
\raggedright
  \item \emph{Say it:} \emph{santān}
  \item \emph{Say it:} \emph{santān}, once more
  \item \emph{Say it:} \emph{temne be santān che, ek dīkro ane ek dīkrī.} — They have two children, a son and a daughter.
  \item \emph{Recall:} say the Gujarati for a respected elder, then the Gujarati for a household, then say \emph{santān} again
\end{itemize}

\begin{tcolorbox}[breakable,colback=teal!4,colframe=teal!35!black,title={Before you move on}]
What does \gu{સંતાન} mean? (\textquotedblleft{}Offspring; a child\textquotedblright{}.) Say it once more.
\end{tcolorbox}
\section[(dialogue)]{Looking back: from Caring for Grandmother to Many Generations}

\label{lesson:GU-R493-review-487-493}



\begin{quote}
Say the Gujarati for nursing care and consoling words.
\end{quote}

\subsection*{Your turn}
Say these aloud:

\begin{itemize}
\raggedright
  \item nursing care, consoling words, a massage, to massage by pressing, to cover someone with a blanket
  \item to bathe someone, to feed, to dress someone, to play with a child, to give someone a drink
  \item faring well, relief from pain, a prayer for someone, to recall, sympathy
  \item a father-in-law, a husband's younger brother, a husband's sister, a husband's elder brother's wife, married life
  \item a son-in-law, a wife's brother, a sister's husband, the father-in-law of one's child, the mother-in-law of one's child
  \item a married woman's parents' home, the in-laws' home, a sister's child, a grandson, a cousin
  \item a generation, an ancestor, a respected elder, a household, offspring
\end{itemize}

\begin{tcolorbox}[breakable,colback=teal!4,colframe=teal!35!black,title={Before you move on}]
Nursing care? (\textbf{\gu{ચાકરી}}, \emph{cākrī}.) Consoling words? (\textbf{\gu{દિલાસો}}, \emph{dilāso}.) A massage? (\textbf{\gu{માલિશ}}, \emph{māliś}.) To massage by pressing? (\textbf{\gu{દાબવું}}, \emph{dābvũ}.) To cover someone with a blanket? (\textbf{\gu{ઓઢાડવું}}, \emph{oḍhāḍvũ}.) To bathe someone? (\textbf{\gu{નવડાવવું}}, \emph{navḍāvvũ}.) To feed? (\textbf{\gu{ખવડાવવું}}, \emph{khavḍāvvũ}.) To dress someone? (\textbf{\gu{પહેરાવવું}}, \emph{paherāvvũ}.) To play with a child? (\textbf{\gu{રમાડવું}}, \emph{ramāḍvũ}.) To give someone a drink? (\textbf{\gu{પીવડાવવું}}, \emph{pīvḍāvvũ}.) Faring well? (\textbf{\gu{કુશળ}}, \emph{kuśaḷ}.) Relief from pain? (\textbf{\gu{શાતા}}, \emph{śātā}.) A prayer for someone? (\textbf{\gu{દુઆ}}, \emph{duā}.) To recall? (\textbf{\gu{સંભારવું}}, \emph{sambhārvũ}.) Sympathy? (\textbf{\gu{હમદર્દી}}, \emph{hamdardī}.) A father-in-law? (\textbf{\gu{સસરો}}, \emph{sasro}.) A husband's younger brother? (\textbf{\gu{દિયર}}, \emph{diyar}.) A husband's sister? (\textbf{\gu{નણંદ}}, \emph{naṇand}.) A husband's elder brother's wife? (\textbf{\gu{જેઠાણી}}, \emph{jeṭhāṇī}.) Married life? (\textbf{\gu{ઘરસંસાર}}, \emph{gharsansār}.) A son-in-law? (\textbf{\gu{જમાઈ}}, \emph{jamāī}.) A wife's brother? (\textbf{\gu{સાળો}}, \emph{sāḷo}.) A sister's husband? (\textbf{\gu{બનેવી}}, \emph{banevī}.) The father-in-law of one's child? (\textbf{\gu{વેવાઈ}}, \emph{vevāī}.) The mother-in-law of one's child? (\textbf{\gu{વેવાણ}}, \emph{vevāṇ}.) A married woman's parents' home? (\textbf{\gu{પિયર}}, \emph{piyar}.) The in-laws' home? (\textbf{\gu{સાસરું}}, \emph{sāsrũ}.) A sister's child? (\textbf{\gu{ભાણેજ}}, \emph{bhāṇej}.) A grandson? (\textbf{\gu{પૌત્ર}}, \emph{pautra}.) A cousin? (\textbf{\gu{પિતરાઈ}}, \emph{pitrāī}.) A generation? (\textbf{\gu{પેઢી}}, \emph{peḍhī}.) An ancestor? (\textbf{\gu{પૂર્વજ}}, \emph{pūrvaj}.) A respected elder? (\textbf{\gu{બુઝુર્ગ}}, \emph{bujhurg}.) A household? (\textbf{\gu{પરિવાર}}, \emph{parivār}.) Offspring? (\textbf{\gu{સંતાન}}, \emph{santān}.)
\end{tcolorbox}
